\section{Discussion}
\label{sec:discussion}

\subsection{Interpretation of Findings}

The results provide evidence that China's 2017 climate-adaptive city
pilot programme enhanced urban ecological resilience, with the effect
operating primarily through accelerated post-shock vegetation recovery
rather than through improved resistance during extreme events. This
finding aligns with the ecological resilience literature, which
emphasises recovery speed as a more tractable policy target than
real-time resistance \citep{smith2023global,li2023widespread}. The
significant XGBoost estimate ($\theta = 0.040$, $p < 0.05$) and the
highly significant recovery-channel estimate ($\theta = 0.081$,
$p < 0.01$) represent, to our knowledge, the first causal evidence
linking climate-adaptive urban policy to measurable improvements in
satellite-derived ecological resilience.

The positive but non-significant Random Forest estimate for CSEE
($\theta = 0.015$, $p \approx 0.32$) warrants discussion. The
difference between learners reflects the well-known tendency of tree
ensembles to underfit low-variance treatment signals, as the CSEE
index's CR component has a standard deviation of only 0.089. XGBoost's
gradient-boosting architecture, which sequentially refines residual
fitting, is better suited to detecting the policy's effect on the
high-variance RC component. The consistency of direction across all
four learners---all produce $\theta > 0$ for CSEE---and the convergence
of the post-2010 and post-2008 subsample estimates to statistical
significance under Random Forest ($\theta = 0.037$ and $\theta = 0.033$,
both $p < 0.05$) suggest that the RF full-sample estimate is
underpowered rather than null.

\subsection{Policy Selection Endogeneity}

An important consideration is whether pilot cities were selected based
on pre-existing environmental trajectories. Our parallel-trends tests
provide reassurance: all five outcome variables pass the joint Wald
$F$-test ($p > 0.21$), indicating no systematic pre-treatment
divergence between pilot and control cities. Furthermore, the
failure of RSEI parallel trends under the combined sponge-city and
climate-adaptive specification ($F = 6.96$, $p < 0.001$) and its
recovery under the single-cohort specification ($F = 1.42$, $p = 0.215$)
demonstrates that the sponge-city cities, not the climate-adaptive
cities, drove the earlier violation. This finding underscores the
importance of policy-specific identification strategies when multiple
overlapping pilot programmes exist.

\subsection{Alternative Outcome: Land Surface Temperature}

As a supplementary analysis, we constructed five land surface
temperature (LST) outcomes from MODIS MOD11A2 data for 338 cities.
While the DML estimates for daytime and nighttime LST are
statistically significant ($\theta_{\mathrm{day}} = 0.35$,
$p < 0.01$), the parallel-trends test for daytime LST fails
($F = 3.19$, $p = 0.008$), indicating that pilot and control cities
had divergent pre-treatment LST trajectories. This violation suggests
that pilot cities may have been selected partly because they face
greater warming pressure---an instance of policy selection
endogeneity that does not affect the CSEE/RSEI outcomes. The LST
results therefore serve as a cautionary illustration of how
selection on outcome trends can produce spurious significance,
reinforcing the importance of parallel-trends validation for any
candidate outcome variable.

\subsection{Heterogeneity and Policy Implications}

The geographic heterogeneity---with eastern cities showing the
strongest effect---is consistent with the hypothesis that policy
implementation effectiveness depends on institutional capacity and
fiscal resources \citep{lu2024spatiotemporal,zhao2025coevolution}.
Eastern China's more developed municipal governments may be better
positioned to execute the pilot programme's ecological restoration
mandates, from procurement of green infrastructure to enforcement of
land-use regulations. The larger effect in smaller cities suggests
that compact urban footprints may facilitate more rapid ecological
recovery, as restoration interventions cover a greater proportion of
the urban ecosystem.

These heterogeneity patterns carry direct policy implications. First,
scaling the climate-adaptive city programme nationally should
prioritise eastern and smaller cities where returns are highest, while
providing additional implementation support to central and western
regions. Second, the recovery-dominated mechanism suggests that policy
design should emphasise post-disturbance restoration
infrastructure---such as rapid replanting protocols, irrigation
systems for drought recovery, and soil remediation---over real-time
protective measures. Third, the stronger effects in the post-2010
subsample indicate that policy benefits accumulate over time,
supporting sustained rather than short-term programme commitments.

\subsection{Limitations}

Several limitations should be noted. First, the CSEE index relies on
NDVI as the sole ecological indicator; while NDVI is the most widely
used satellite vegetation measure, it does not capture biodiversity,
soil health, or below-ground processes. Future research could extend
the index using enhanced vegetation index (EVI), solar-induced
fluorescence (SIF), or species richness data. Second, the treatment
effect is identified through a single policy cohort (2017), which
limits our ability to examine treatment-effect dynamics across
multiple cohorts. Third, the 27 treated cities represent a relatively
small treatment group, which constrains statistical power for
subgroup analysis. Fourth, while we control for a rich set of
observables, unobserved time-varying confounders---such as concurrent
environmental regulations---could bias estimates if they correlate
with both pilot status and ecological outcomes. The placebo test
and PSM results provide partial reassurance, but cannot fully rule
out this concern.
